\begin{proposition}[Proved: the plumbing equation for the corrected inverse]
\label{A:cap-plumbing-equation}
The inverse on the restricted propagator domain in
Theorem~\ref{A:graph-radius-domain} is a plumbing in the canonical cap
coordinates, which themselves depend on the recovered length vector.
If $z_\pm$ are coordinates at the two node branches and
$f_\pm:\D\to D_\pm$ are normalized so that
$(z_\pm\circ f_\pm)'(0)=r_\pm>0$, put
$Z_\pm=r_\pm f_\pm^{-1}$. Its seam identification is exactly
\[
                       Z_+Z_-=q,\qquad
             q=r_+r_-e^{i\theta}.
\]
This identity does not identify the radius inverse with plumbing in
one fixed pair of coordinate germs for all $q$.
\end{proposition}

\begin{proof}
The canonical cap boundary is analytic, so its normalized conformal
coordinate extends univalently to a collar across the boundary.
On that boundary the toric identification is
$f_-^{-1}=e^{i\theta}/f_+^{-1}$, proving the displayed equation.
Take slightly enlarged coordinate discs of radii $R_\pm>r_\pm$.
They have $R_+R_->|q|$, so the transition $Z_-=q/Z_+$ identifies
overlapping annuli. Its restriction to the seam is the specified
identification; its holomorphic extension gives the same conformal
welding by uniqueness. The derivatives of $Z_\pm$ relative to
$z_\pm$ at the origins are one, so the coefficient is the actual
tensor coefficient used in Lemma~\ref{A:tensor}.

The inverse radius theorem first recovers the lengths, and only then
these cap maps. Their dependence is real analytic on each face.
Higher jets of $Z_\pm$ therefore move with $q$. The tensor supplies
the first derivatives and product, whereas a fixed-coordinate
plumbing also specifies those higher jets. Equality of the equations
in their respective moving coordinates is consequently the precise
comparison that the construction proves.
\end{proof}

\begin{proposition}[Known: the applicable horocyclic and metric comparisons]
\label{A:plumbing-sources}
Kra's horocyclic construction and Wolpert's hyperbolic plumbing
are valid comparisons after their coordinate choices and domain
restrictions are retained. They do not parametrize arbitrary
nonzero normal tensors by the manuscript's canonical capping map.
\end{proposition}

\begin{proof}
For a maximal pants graph with $d=3g-3+n$, \citet[Theorems~1
and~3, pp.~501--503; \S\S2--3]{Kra1990} uses cusp coordinates and
the equation $zw=t$. His tame construction contains the polydisc
$|t|_{\max}<e^{-2\pi}$. The full domain $D_\Gamma$ is a domain of
holomorphy satisfying
\[
  (\D_{e^{-2\pi}})^d\subset D_\Gamma
                    \subset(\D_{e^{-\pi/2}})^d,
\]
and its map to ordinary moduli can have infinitely many preimages.
The logarithmic coordinates on the fully marked Teichmuller space
have Dehn twists $\tau_j\mapsto\tau_j+2$.
Thus his domain and marking quotient must both be kept.
The present inverse instead chooses its discs by radius inversion.
Proposition~\ref{A:cap-plumbing-equation} proves the shared local
equation without identifying these different coordinate choices.

For the hyperbolic comparison, use normalized cusp coordinates
$z=e^{2\pi i\zeta}$ and $w$ on the limiting components. On
$|t|<|z|<1$, the complete annular metric is
\[
 ds_t^2=\left(
 \frac{\pi}{\log|t|}
 \csc\frac{\pi\log|z|}{\log|t|}
 \left|\frac{dz}{z}\right|\right)^2.
\]
Integrating on $|z|=|t|^{1/2}$ gives core length
$2\pi^2/(-\log|t|)$. Wolpert's curvature correction transfers this
model to the plumbed surface, with normalized parameter
$\tau_j$ including the first derivatives of the chosen coordinates:
\citet[Expansion~4.2, p.~445, and Example~4.3, p.~446]{Wolpert1990}.
In particular, for one pinched curve and fixed remaining data,
\[
 \ell(t)=\frac{2\pi^2}{-\log|t|}
                   +O((-\log|t|)^{-2}).
\]
The coordinate scaling changes the logarithm by a bounded constant,
which explains this form of the error.
The punctured-surface version and its hypotheses are stated in
\citet[\S2, Theorems~4 and~6]{ObitsuWolpert2008}: the local
parameters are real analytic in the normalization variables and
holomorphic in $t$ with those variables fixed. Here their metric
interpolation called grafting is the construction in that paper,
distinct from attaching the infinite flat cap cylinders.

Finally Hubbard--Koch's local chart theorem, applied in
Proposition~\ref{A:plumbing}, identifies these plumbing germs with
the analytic Deligne--Mumford deformation problem. The first-order
transition $t'=ab\,t$ proves the common normal-line interpretation.
The independent unrestricted-vector counterexample in
Theorem~\ref{A:04} excludes any global inverse of the kind proposed
in L2, regardless of which of these local coordinate theorems is
used.
\end{proof}

\begin{proposition}[Known: the proved local Earle--Marden comparison]
\label{A:EarleMarden-local}
On a neighbourhood of a fixed nodal surface, the Earle--Marden
Kleinian plumbing family and the Hubbard--Koch family give the same
analytic marked deformation germ. This comparison follows from the
proved Appendix of \citet[pp.~41--42 in the opened author version]
{HubbardKoch2014}. It is a local comparison with a specified
Dehn-twist quotient.
\end{proposition}

\begin{proof}
We spell out the identification used here. Normalize each peripheral
parabolic of the nodal Kleinian group to $\zeta\mapsto\zeta+1$.
The added translation by $\tau_e$, with imaginary part exceeding a
fixed collar bound, opens that node. Passing to
$t_e=e^{2\pi i\tau_e}$ removes exactly the integral Dehn-twist
ambiguity. With the transition $\zeta_+=\zeta_-+\tau_e$, set
$z_e=e^{2\pi i\zeta_+}$ and $w_e=e^{-2\pi i\zeta_-}$.
Then $z_ew_e=e^{2\pi i\tau_e}=t_e$ exactly; in particular the
smoothing exponent is one, after these fixed normalizations.
Together with the parameters $u$ of the normalization components,
this constructs the $\Gamma$-marked analytic stable family on
$U\times\D_\rho^E$ described in the cited Appendix.

By Hubbard--Koch's Theorem~10.1 its universal property supplies
an analytic classifying map to their $\mathcal Q_\Gamma$.
At the nodal point, the $u$-directions are the same normalization
deformations on both sides; each $t_e$ opens just its own node
and has its nonzero normal tangent coefficient in
$T_{p_e^+}\otimes T_{p_e^-}$. Thus the derivative is an
isomorphism, by the independent smoothing directions and
normal-line calculation in Proposition~\ref{A:plumbing}.
The analytic inverse function theorem identifies the germs, and
the classifying property identifies their families. This is the
comparison proved by the Appendix, not a claim that the
neighbourhood can be chosen to be the whole tensor bundle.

The bibliography separately marks the original 2012
Earle--Marden chapter as not opened. We have verified the
mathematical comparison in the opened primary proof just
specified. The prompt does not identify the 2012 chapter as the
required source: the opened Appendix is itself a primary proof of the
requested construction comparison. The unavailable chapter remains a
bibliographic coverage limitation, and no statement is accepted on its
unopened authority.
\end{proof}

\begin{example}[Refuted: arbitrary fixed plumbing is everywhere a local chart]
\label{A:hinich-explicit}
Even a holomorphic plumbing of two three-pointed spheres with
disjoint coordinate discs can have zero derivative at a nonzero
parameter. We give exact rational data for the mechanism of
\citet[\S\S2.1--2.3]{Hinich2010}, rather than leaving its
parameters existential.
\end{example}

\begin{proof}
Set
\[
 R=\frac32,\quad B=\frac34,\quad C=10,\quad
 A=\frac{420}{569},\quad\varepsilon=\frac1{100},\quad
 \phi(\zeta)=(\zeta-R)^2.
\]
The map $\phi$ is injective on $|\zeta|<1+\varepsilon$:
two different points with the same image would sum to $2R=3$,
whereas their sum has modulus below $2(1+\varepsilon)<3$.
Its derivative has no zero there. Put
$a=\phi(A)$, $b=\phi(B)$, $c=\phi(C)$. The distinct marks
$a,b$ lie on the bounded side of
$\phi(\{1-\varepsilon<|\zeta|<1+\varepsilon\})$.
The marks $c,\infty$ lie on the other side, since both roots
$C,2R-C$ of $\phi(\zeta)=c$ have modulus above $1+\varepsilon$.

We specify the two caps to avoid any orientation ambiguity.
Attach a disc with centre $\zeta=\infty$, using the region
$|\zeta|>1-\varepsilon$, to the bounded-side component and its
annular collar. This makes the first sphere, marked by $a,b$
and the new centre. Attach the disc $|\zeta|<1+\varepsilon$,
with centre zero, to the unbounded-side component and its collar.
This makes the second sphere, marked by $c,\infty$ and its
new centre. The overlaps are identified by $\phi$.
Each compactified surface is a topological sphere with a
holomorphic atlas, hence is a Riemann sphere. The respective
node coordinates are
\[
             x=(1-\varepsilon)/\zeta,\qquad
             y=\zeta/(1+\varepsilon).
\]
Their unit discs contain none of the other markings, and lie
on distinct components, so the disjointness condition holds.
Plumbing $xy=t$ at
$t_0=(1-\varepsilon)/(1+\varepsilon)=99/101$
recovers the sphere with marks $a,b,c,\infty$.

The Kodaira--Spencer cocycle of a nonzero $t$-variation is a
nonzero constant times $\zeta\partial_\zeta$ on this annulus.
The cotangent space to the labelled four-point moduli at this
sphere is one-dimensional, with basis
$Q=dw^2/((w-a)(w-b)(w-c))$. This follows either from Serre
duality or from $K_{\PP^1}^{\otimes2}(a+b+c+\infty)$ having
degree zero. The cocycle pairs with $Q$ by the contour integral
of
\[
 \frac{\zeta\phi'(\zeta)^2\,d\zeta}
      {(\phi(\zeta)-a)(\phi(\zeta)-b)(\phi(\zeta)-c)}.
\]
The only poles inside the unit circle are $A,B$. The sum of
their residues is
\[
 \frac{2}{(A-B)(A+B-2R)}
 \left[
 \frac{A(A-R)}{(A-C)(A+C-2R)}
 -
 \frac{B(B-R)}{(B-C)(B+C-2R)}
 \right]=0.
\]
Indeed solving this cancellation, after cancelling $A-B$,
gives $A=C(R-B)(C-2R)/(C^2-2RC+RB)=420/569$.
Thus the pairing with the whole cotangent space vanishes,
and the tangent map is zero. The family is a local chart
near its nodal parameter by Proposition~\ref{A:plumbing},
so this is a nonconstant family with a critical point.
The exact-arithmetic program
\texttt{code/A\_plumbing\_certificate.py} verifies all ten
finite parameter and residue checks, including the locations
of all six distinct poles. This example does not contradict
the local theorem near $t=0$ or the everywhere invertible
real radius map, which uses moving canonical cap coordinates.
% Suggested insertion in the proof of A:hinich-explicit.
% Independent audit by /root/algebra_c, 19 November2026.
\paragraph{An independent explicit cross-ratio verification.}
For these particular coordinates one can compute the entire plumbing family directly, without using the cotangent pairing. The first capped component has a global coordinate $\zeta_1$: on the original bounded component it is the inverse of the injective map $w=\phi(\zeta_1)$, and it is the exterior coordinate on its attached cap. Its three marks are $A,B,\infty$. The second capped component has the original global $w$-coordinate: on its attached interior cap use $w=\phi(\zeta_2)$. Its three marks are $c,\infty,R^2$. The two node coordinates are therefore exactly
\[
 x=(1-\varepsilon)/\zeta_1,
 \qquad y=\zeta_2/(1+\varepsilon).
\]

Fix any $t$ with $0<|t|<1$, and put $\alpha=t/t_0$. The first retained plumbing piece is
\[
 |\zeta_1|<(1-\varepsilon)/|t|,
\]
and the second retained piece is the complement in the $w$-sphere of
\[
 \phi\bigl(\{|\zeta_2|\le(1+\varepsilon)|t|\}\bigr).
\]
On the first piece define $w_t=\phi(\alpha\zeta_1)$; on the second use its existing $w$-coordinate. The plumbing equation $xy=t$ is equivalent to $\zeta_2=\alpha\zeta_1$, so these functions agree on the overlap. Because
\[
 |\alpha|\,(1-\varepsilon)/|t|
       =(1-\varepsilon)/t_0=1+\varepsilon<R,
\]
and $\phi$ is injective on $|\zeta|<1+\varepsilon$, the first piece maps biholomorphically to $\phi(\{|\zeta|<1+\varepsilon\})$. Its overlap with the second piece is exactly the image of
\[
 (1+\varepsilon)|t|<|\zeta_2|<1+\varepsilon.
\]
The two images cover the $w$-sphere. Thus $w_t$ is a global conformal coordinate on the plumbed sphere for every $0<|t|<1$, not merely near $t_0$.

The four labelled marks in this coordinate are
\[
 a(t)=\phi((t/t_0)A),\quad
 b(t)=\phi((t/t_0)B),\quad c=\phi(C),\quad\infty.
\]
They stay distinct throughout this punctured unit disc: the first two lie in the injective $\phi$-disc, while $c,\infty$ are outside it. A cross-ratio coordinate on labelled $M_{0,4}$ is therefore the explicit rational function
\[
 \Lambda(t)=\frac{\phi((t/t_0)A)-c}{\phi((t/t_0)B)-c}.
\]
Substituting the specified rational constants and differentiating gives
\[
 \Lambda'(t)=\frac{201579840}{323761}\,
 \frac{(101t-99)(101t+12518)}
 {(101t-1320)^2(101t+924)^2}.
\]
In particular
\begin{align*}
 \Lambda(t_0)&=\frac{323680}{323761},
 &\Lambda'(t_0)&=0,\\
 \Lambda'(0)&=-\frac{909}{1752520}\ne0,
 &\Lambda''(t_0)&=\frac{22850240}{44933817907}\ne0.
\end{align*}
Thus the nonzero parameter $t_0$ is a simple critical point, while the family is a local smoothing coordinate at the nodal parameter $0$. This verifies the geometric conclusion independently of the Kodaira--Spencer computation. The exact certificate \texttt{code/C\_plumbing\_crossratio.py}, with output \texttt{logs/C\_plumbing\_crossratio.json}, also checks that all zeros and poles of $\Lambda$ lie outside the closed unit disc, and that $\Lambda=1$ occurs there only at $t=0$. Its eleven assertions use exact rational functions.

\end{proof}
